% =============================================================================
% CONCLUSION — très courte, après le Bloc 4
% =============================================================================
\section{Conclusion}
\begin{frame}{Quatre blocs, une progression professionnelle}
\context{CONCLUSION}{Synthèse}
\begin{columns}[T]
\begin{column}{0.48\linewidth}
\begin{block}{Ce que les huit projets démontrent}
\begin{tightitems}
  \item \strong{Planifier} : cadrer, décider, mesurer et réajuster.
  \item \strong{Concevoir} : choisir le bon niveau d'abstraction et développer une solution adaptée.
  \item \strong{Industrialiser} : versionner, tester, déployer, observer et documenter.
  \item \strong{Piloter} : organiser un collectif, accompagner et donner du feedback.
\end{tightitems}
\end{block}
\end{column}
\begin{column}{0.48\linewidth}
\begin{block}{Mon principal apprentissage}
Je suis passé d'une logique centrée sur « faire fonctionner le code » à une logique où la qualité d'une solution dépend aussi du \textbf{contexte métier}, du \textbf{risque}, de la \textbf{production}, de la \textbf{documentation} et des \textbf{personnes qui doivent la faire vivre}.
\end{block}
\vspace{0.5em}
\centering
{\Large\color{Navy}\textbf{Questions}}
\end{column}
\end{columns}
\end{frame}
